\documentclass[aps,prx,article,twocolumn,amsmath,amssymb,superscriptaddress,longbibliography]{revtex4-2}
\usepackage{graphicx}
\usepackage{dcolumn}
\usepackage{bm}
\usepackage{amssymb}
\usepackage{amsmath}
\usepackage{mathrsfs}
\usepackage{epigraph}
\usepackage{braket}
\usepackage{gensymb}
\usepackage{lmodern}
\usepackage{tipa}
\usepackage{bbold}
\usepackage{esint}
\usepackage{mathdots}
\usepackage{xcolor}
\usepackage{appendix}
\usepackage{multirow}
\usepackage{float}
\usepackage{physics}
\usepackage[normalem]{ulem}
\usepackage{times}
\usepackage{comment}
\usepackage[colorlinks=true, linkcolor=blue, urlcolor=blue, citecolor=blue]{hyperref}

\begin{document}

\title{Unconventional Hysteresis Induced by Coupling of Stripe-like Charge Orders and Lattice Anharmonicity}
\author{Haoran Yan}
\affiliation{Department of Chemistry, Emory University, Atlanta, GA 30322, USA}
\author{Ziyan Zhu}
\affiliation{Department of Physics, Boston College, Chestnut Hill, MA, USA}
\author{Yao Wang}
\email[Correspondence should be addressed to \href{mailto:yao.wang@emory.edu}{yao.wang@emory.edu}]{}
\affiliation{Department of Chemistry, Emory University, Atlanta, GA 30322, USA}
\date{\today}

\begin{abstract}
(to be changed) Bilayer graphene provides a minimal platform in which layer polarization, electrostatic displacement fields, and moiré superlattice potentials can be combined to engineer interaction-driven polar electronic states. Motivated by recent observations of ferroelectric-like hysteresis in graphene-based moiré systems, we develop a theoretical framework for excitonic ferroelectricity in bilayer graphene with a layer-contrasting moiré structure. The key idea is that the moiré potential acts differently on the two graphene layers, generating low-energy states with spatially modulated layer character. Coulomb interactions between these layer-polarized electron and hole states can then drive an interlayer excitonic instability. Once established, the excitonic order breaks inversion symmetry, locks to the layer-contrasting moiré background, and produces an electrically switchable polarization. This framework connects excitonic coherence, moiré band engineering, and ferroelectric response in bilayer graphene, and provides a route for understanding hysteretic transport phenomena without invoking a purely structural ferroelectric transition.
\end{abstract}

\maketitle

\section{Introduction}

\textcolor{blue}{\itshape \textbf{Paragraph 1}
Flat-band two-dimensional systems provide a tunable platform for correlated and topological quantum phases. BLG as an example. insulating states, quantum Hall ferromagnetism, superconductivity, and electrically tunable layer-polarized or polar electronic states.}

Flat-band two-dimensional systems provide an ideal platform for realizing interaction-driven quantum phases. Owing to the reduced kinetic energy and enhanced interactions, these materials host a wide range of correlated phenomena. Among these materials, bilayer graphene systems (BLG) are particularly attractive because its electronic structure can be continuously tuned by carrier density and external electric fields. As a result, a variety of electrically controllable ordered states have been observed, including insulating states, quantum Hall ferromagnetism, superconductivity, and electronic symmetry-breaking orders.

\textcolor{blue}{\itshape \textbf{Paragraph 2}
Discovered charge-ordered states in moiré BLG. The charge order or stripe-like states may coexist with and even enhance superconductivity, indicating a potential connection between electronic ordering and unconventional pairing. Earlier SNOM experiments indicated signatures of phonons, and earlier theory suggested phonon-driven superconductivity (Das Sarma). These observations indicate possible connections among charge order, phonons, and superconductivity.}

Beyond these magnetism and topological phases, increasing experimental attention has recently focused on charge ordering and stripe-like electronic states in BLG. Transport measurements have revealed pronounced electronic anisotropy and symmetry breaking near regions of enhanced density of states (DOS). At the same time, spectroscopic studies have identified substantial electron--phonon coupling in BLG, indicating that lattice degrees of freedom may actively participate in the formation and stabilization of these ordered states. These observations indicate that charge order, electronic anisotropy, and lattice distortions are intimately connected manifestations of the same underlying instability.

\textcolor{blue}{\itshape \textbf{Paragraph 3}
After the charge-order is introduced, discuss about the related hysteresis. Mention the phenomenons found in rachite experiments.
}

Recent experiments have further revealed doping-driven hysteretic behavior accompanying the emergence of electronic orders and symmetry-broken states. The hysteresis develops near charge-ordering regimes and is highly sensitive to details of the band structure, indicating that it is not simply a consequence of a bistable structural configuration or conventional sliding ferroelectricity. Rather, it points to a mechanism in which electronic degrees of freedom actively participate in the switching process. Similar highly reproducible and programmable memory states have been observed in the quasi-one-dimensional quantum spin Hall insulator TaIrTe$_4$, where scanning-probe and Raman measurements attributed the bistability to strong electron--phonon coupling,\cite{tang2026bistable}. These observations motivate a unified picture in which hysteresis in correlated quantum materials emerges from the coupled dynamics of electronic ordering and anharmonic lattice degrees of freedom.


\textcolor{blue}{\itshape \textbf{Paragraph 4}
State our proposed mechanism. }

Based on these observations, we introduce a minimal phenomenological model to understand the connection between hysteresis and electronic charge order in BLG. The same framework can also be naturally extended to other flat-band materials with pronounced lattice anharmonicity. Our central proposal is that, near the flat-band VHS regime, the enhanced density of states and Fermi-surface nesting drive the formation of a strong stripe-like electronic charge order, which breaks both translational and rotational symmetries. Through electron--lattice coupling, this electronic order can further induce a substantial lattice distortion toward a lower-symmetry structure, with a superlattice periodicity commensurate with the CDW order. Importantly, the resulting lattice-distorted state can remain metastable even after the electronic order is suppressed when the doping is tuned away from the VHS. This mechanism, distinct from the conventional picture of a simple bistable structural configuration, is consistent with recent observations of multi-level hysteresis and the reduction of conducting carriers in BLG. It therefore provides a unified explanation for the recently observed stripe order, charge redistribution, and hysteretic memory states.





\section{Electronic and Lattice Instabilities}

\textcolor{blue}{\itshape \textbf{Paragraph 5}
Introduce the BLG electronic instability from layer-contrasting moiré minibands. The relevant flat bands have enhanced DOS and strong layer texture, so interactions can drive finite-momentum charge order or an excitonic/layer-polarized density modulation. This order reconstructs the minibands and reduces the density of itinerant carriers visible in transport.}

We now turn to the electronic instability that may underlie the anomalous transport response in moiré BLG. In Bernal BLG aligned with hBN, the moiré potential reshapes the low-energy electronic structure into narrow minibands whose wave functions carry strong layer character. When the moiré environment is layer asymmetric, the two graphene layers experience distinct periodic potentials, so the low-energy states are not only flat but also spatially and layer textured. This combination of enhanced DOS, layer polarization, and moiré form factors makes the system susceptible to interaction-driven ordering instabilities in charge, excitonic, or layer-pseudospin channels. The relevant ordering tendency is therefore controlled by the moiré miniband geometry and by the layer-contrasting character of the Bloch states.

The central consequence of such an electronic order is a reconstruction of the low-energy conducting states. A finite-momentum charge modulation folds the moiré bands and hybridizes states separated by the ordering wavevector, while an excitonic or layer-polarized density order coherently transfers spectral weight between layer-polarized bands. In either case, part of the spectral weight that would otherwise contribute to ordinary Drude transport can become locked into a collective ordered state. The system therefore need not become a full insulator in order to exhibit a reduced density of itinerant carriers. This provides the electronic basis for interpreting the apparent loss of conducting electrons observed in transport as a signature of an underlying ordered state.

To quantify the tendency toward charge ordering, we evaluate the static charge susceptibility within the generalized Lindhard formalism in the non-interacting limit,
\begin{align}\label{eq:susceptibility}
\chi({\vb*{q}}) =
\sum_{\mathbf{k},m,n}|\left\langle u_{m,\mathbf k+\mathbf q} \middle| u_{n,\mathbf k}\right\rangle|^2
\frac{f(\epsilon_{m,\mathbf{k}})-f(\epsilon_{n,\mathbf{k}+{\vb*{q}}})}
{\epsilon_{n,\mathbf{k}+{\vb*{q}}}-\epsilon_{m,\mathbf{k}}+i\eta},
\end{align}
where $m$ and $n$ label moiré minibands, $f(\epsilon)$ is the Fermi--Dirac distribution function, and $\eta$ is a small \textit{ad hoc} broadening parameter introduced to regularize the susceptibility. The form factor $\left\langle u_{m,\mathbf k+\mathbf q} \middle| u_{n,\mathbf k}\right\rangle$ captures the overlap between layer- and moiré-textured Bloch wave functions. It is therefore essential in BLG, where the ordering tendency depends not only on the DOS but also on whether the relevant states occupy the same layer, opposite layers, or different regions of the moiré unit cell.

Based on the susceptibility, we formulate an effective Ginzburg--Landau theory for the dominant electronic order. We denote the order parameter by $\psi_{\vb*{q}}$, where $\vb*{q}$ is selected by the moiré miniband structure and can represent either a charge-density modulation or the charge component of an excitonic/layer-polarized order. The microscopic interactions and wave-function form factors determine the precise ordering channel, but the phenomenological structure is common to these possibilities. The corresponding electronic free energy takes the form
\begin{align}\label{eq:electronFreeEnergy}
\mathcal{F}_{\rm ele} =
a(T,n)\abs{\psi_{\vb*{q}}}^2 + b\abs{\psi_{\vb*{q}}}^4,
\end{align}
with the second-order coefficient
\begin{align}\label{eq:electronCoefficient}
a(T,n) = \frac{1}{V_{\vb*{q}}} - \chi_{\vb*{q}}(T,n),
\end{align}
where $V_{\vb*{q}}$ denotes the effective interaction in the ordering channel after high-energy degrees of freedom are integrated out. When the susceptibility is enhanced by the flat moiré minibands and exceeds $1/V_{\vb*{q}}$, the coefficient $a(T,n)$ becomes negative and the electronic order develops. If the dominant component of $\psi_{\vb*{q}}$ is a charge modulation, the ordered state breaks translational symmetry within the moiré superlattice. If the dominant component is layer-polarized or excitonic, it also breaks inversion or layer-pseudospin symmetry. In both cases, the ordered state can strongly modify the transport carrier density by reconstructing the conducting minibands.

The hysteresis observed in BLG further suggests that this electronic instability is coupled to a slow metastable degree of freedom. A purely electronic order that appears only in a narrow range of carrier density would normally disappear once the susceptibility is reduced. By contrast, a history-dependent response requires a free-energy landscape with multiple local minima. In BLG/hBN devices, natural candidates for such slow coordinates include local moiré registry, interfacial relaxation, layer-asymmetric electrostatic polarization, and long-wavelength lattice distortions of the graphene and hBN layers. We describe these effects phenomenologically by a lattice or moiré distortion coordinate $X_{\vb*{q}}$ that couples to the electronic order at the same wavevector.

The corresponding lattice sector is expected to be anharmonic because the moiré interface samples multiple nearly degenerate local stacking environments. We therefore do not assume that $X_{\vb*{q}}$ is a simple harmonic phonon coordinate. Instead, it represents a collective structural or electrostatic configuration that can be pinned by the substrate and trapped in metastable minima. Once the electronic order forms, electron--lattice coupling can displace this coordinate and imprint the electronic symmetry breaking into the moiré structure. After the carrier density is tuned away from the strongest electronic instability, the structural coordinate may remain trapped, allowing the reconstructed electronic state to persist over a broader doping range. This coupled electronic--lattice picture provides a natural route to hysteresis and to the apparent disappearance of conducting carriers in layer-contrasting moiré BLG.


















\section{}






\bibliography{refs}

\end{document}
